\chapter{実験}
\label{chap:実験}

本章では，提案システムの基盤となる音楽生成モデルRenderBoxの訓練詳細，および提案システムの有効性を検証するために実施した被験者実験について述べる．

\section{RenderBoxの訓練}
\label{sec:RenderBoxの訓練}

\subsection{データセット}
RenderBoxの訓練には，ASAP\cite{asap-dataset}，MusicNet\cite{musicnet-dataset}，GAPS\cite{gaps-dataset}，BachViolin\cite{bach-violine-dataset}，ATEPP\cite{atepp-dataset}，ConEspressione\cite{conespressione-dataset}，Vienna4x22\cite{vienna4x22}を使用した．
各データセットは，訓練用データと推論（評価）用データに9対1の割合でランダムに分割して使用した．
各ステージにおけるデータセットの割り当てを表\ref{tab:datasets}に示す．

\begin{table}[htbp]
\centering
\caption{カリキュラム学習の各ステージで使用したデータセット\cite{RenderBox}}
\label{tab:datasets}
\resizebox{\textwidth}{!}{%
\begin{tabular}{llllll}
\hline
\toprule
\textbf{Datasets by stage} & \textbf{Input MIDI} & \textbf{Target Audio} & \textbf{Inst.} & \textbf{Controls / Prompts} & \textbf{Repertoire / Notes} \\ \midrule
\hline
% ▼▼▼ ここを {7} から {6} に修正しました ▼▼▼
\multicolumn{6}{l}{\textit{Stage 0: Synthesis}} \\
ASAP & Score \& & Synth. Score \&  & Piano & - & Common Practice Period \\
 & Performance & Perf. Audio & & & solo piano works \\
MusicNet & Score \& & Synth. Score \&  & Orchestral & - & Classical music with a large \\
 & Performance & Perf. Audio & & & portion of chamber works \\
GAPS & Score & Synth. Score & Guitar & - & Noted classical guitar works \\ \midrule
\hline
% ▼▼▼ ここを {7} から {6} に修正しました ▼▼▼
\multicolumn{6}{l}{\textit{Stage 1: Speed-controlled synthesis}} \\
ASAP & Score & Synth. Score  & Piano & Tempo & Time-varied score synthesis \\
GAPS & Score & Synth. Score & Guitar & Tempo & Time-varied score synthesis \\
MusicNet & Score & Synth. Score & Orchestral & Tempo & Time-varied score synthesis \\ \midrule
\hline
% ▼▼▼ ここを {7} から {6} に修正しました ▼▼▼
\multicolumn{6}{l}{\textit{Stage 2: Expressive Performance}} \\
ASAP & Score & Perf. Audio & Piano & Expressive & Aligned perf. segments \\
GAPS & Score & Perf. Audio & Guitar & Expressive & Aligned perf. segments \\
MusicNet & Score & Perf. Audio & Orchestral & Expressive & Aligned perf. segments \\
BachViolin & Score & Perf. Audio & Violin & Expressive & Bach violin sonata \\ \midrule
\hline
% ▼▼▼ ここを {7} から {6} に修正しました ▼▼▼
\multicolumn{6}{l}{\textit{Stage 4: Style-directioned Performance}} \\
ATEPP-subset & Score & Perf. Audio & Piano & Performer labels & Standard piano repertoire \\
 & & & & &  played by 40+ pianists \\
Con Espressione & Score & Perf. Audio & Piano & Human-annotated & 45 perf. of 9 excerpts annot. \\
 & & & & expression labels & by 179 participants \\
Vienna 4x22 & Score & Perf. Audio & Piano & Performer labels & 4 short pieces played by 22 \\
 & & & & & pianists each \\ \bottomrule
 \hline
\end{tabular}%
}
\end{table}

\subsection{カリキュラム学習}
\label{subsec:カリキュラム学習}

本研究では，RenderBoxの学習においてカリキュラム学習を採用した．学習は段階的にタスクの難易度や条件を変更しながら進められる．
元のRenderBoxの論文では，意図的に演奏ミスを含ませてモデルのロバスト性を向上させるStage 3が存在するが，本研究では「理想的な表現生成」に焦点を当てるため，このStage 3は省略し，Stage 0, 1, 2, 4の順で学習を行った．

学習環境として，GPUにはNVIDIA A100 (80GB)を2〜3枚使用した．学習時のハイパーパラメータについては，バッチサイズを108〜112に設定し，計算資源の効率的な利用を図った．

\subsection{性能評価指標}
\label{subsec:性能評価指標}

本研究では，生成された演奏の品質，楽曲の整合性，およびテキスト指示への追従性を評価するために，Evansら\cite{stable-audio-open}によって確立され，stable-audio-metricsに実装されている指標を採用した．具体的には以下の指標を用いる．

\begin{itemize}
    \item \textbf{FAD (Fréchet Audio Distance):}
    生成された音声の品質と多様性を評価する指標である．Cramerら\cite{FAD}によるOpenL3埋め込み表現を用いて，生成音声の分布と正解データ（Ground Truth）の分布間の距離を算出する．なお，本実験におけるFADの算出にあたっては，生成された演奏群と，その楽曲が属するデータセットのテスト用データセット（Ground Truth）との間で比較を行った．
    
    \item \textbf{CLAP Score:}
    テキストプロンプトと生成音声の意味的な整合性を評価する指標である．Wuら\cite{clap}によるLAION-CLAP空間における，入力テキストプロンプトと出力音声の埋め込みベクトル間の距離として計算される．
    
    \item \textbf{Chroma Similarity (Pitch Accuracy):}
    演奏生成タスクにおいては，正しい楽曲（楽譜通りの音高）が演奏されていることが不可欠である．そこで，ピッチ方向の正確性を評価するために，MusicGen-Melody\cite{copet2023}の評価手法に着想を得たChroma類似度を用いる．
    生成音声と正解音声は必ずしも時間的に整合（タイムアラインメント）していないため，まずクロマグラムに対して動的時間伸縮（DTW: Dynamic Time Warping）を適用し，大まかなフレーム対応を推定する．このアライメントに基づき，対応するフレーム間のコサイン類似度を計算する．さらに，アライメントのコスト（DTWコスト）をペナルティ項として組み込み，重み係数 $\lambda = 10^{-3}$ でスケーリングして最終的なスコアとする．
    
    \item \textbf{Tempo Deviation:}
    生成された演奏のテンポが，プロンプトで指定されたテンポ（Table \ref{tab:datasets}参照）の範囲内に収まっているかを評価する．
    \texttt{madmom}ライブラリを用いて生成音声のテンポを推定し，プロンプトによって調整された楽譜上の目標テンポに対する差異（逸脱度）をパーセンテージで算出する．
\end{itemize}

\section{被験者実験}
\label{sec:被験者実験}

提案システムにおける「Human-in-the-Loop」および「LLMによる仮説生成」の有効性を検証するため，被験者実験を行った．

\subsection{被験者実験の概要}
\label{subsec:被験者実験の概要}

被験者は男性4名，女性4名の計8名である．被験者の音楽的背景（楽器経験年数，普段聴く音楽ジャンル等）については，実験後のアンケートにより収集した．被験者の属性詳細を図\ref{fig:participants_demographics}に示す（または後述の表\ref{tab:questionnaire_items}のD項にて収集）．

% 図を入れる場合
% \begin{figure}[htbp]
% \centering
% \includegraphics[width=0.8\textwidth]{figures/demographics.png}
% \caption{被験者の音楽的バックグラウンド}
% \label{fig:participants_demographics}
% \end{figure}

本実験では，システムが提示する演奏に対するフィードバックを通じて，ユーザーが理想とする演奏表現を探索するプロセスを評価対象とする．

\subsection{被験者実験の手順}
\label{subsec:被験者実験の手順}

実験は，共通の導入フェーズである「実験0」を行った後，比較対象となる2つの条件「実験1（手動プロンプト）」と「実験2（LLMプロンプト生成）」を順不同（またはカウンターバランスをとって）実施した．具体的な手順は以下の通りである．

\subsubsection*{実験0：初期ループ（共通）}
まず，システムの使用感に慣れるため，以下の手順で最初の演奏生成を行う．
\begin{enumerate}
    \item ユーザーは課題曲のScoreに基づく機械的な演奏を聴取する．
    \item その曲に対する「理想の感性的な表現」を言語化し，システムに入力する（図\ref{fig:proposed_system}の(1)(2)）．
    \item 入力されたテキストを直接プロンプトとしてRenderBoxに入力し，音楽を生成する（図\ref{fig:proposed_system}の(4)(5)）．
    \item ユーザーは生成された演奏を聴き，フィードバック（コメント）を行う（図\ref{fig:proposed_system}の(6)(1)）．このとき生成された演奏を，以降の実験における「基準演奏」とする．
\end{enumerate}

\subsubsection*{実験1：人が直接プロンプトを入力（ベースライン）}
ユーザー自身が表現を言語化し，試行錯誤する条件である．
\begin{enumerate}
    \item ユーザーは基準演奏に対し，改善したい点や新たなイメージをプロンプトとして直接入力する（図\ref{fig:proposed_system}の(2)）．
    \item 入力されたプロンプトをもとに，RenderBoxが音楽を生成する（図\ref{fig:proposed_system}の(5)）．
    \item ユーザーは生成演奏に対し，フィードバック（コメントと満足度ランキング）を行う（図\ref{fig:proposed_system}の(6)(1)）．
    \item ユーザーが演奏に満足するか，または規定のループ回数上限に達するまで，手順1〜3を繰り返す．
\end{enumerate}

\subsubsection*{実験2：LLMがプロンプトを生成（提案手法）}
LLMがユーザーのフィードバックを解釈し，多様な表現仮説を提案する条件である．
\begin{enumerate}
    \item ユーザーのフィードバック（コメントとランキング）をLLMに入力する（図\ref{fig:proposed_system}の(2)）．
    \item LLMはフィードバックに基づき，以下の3種類の仮説（プロンプト）を立案する（図\ref{fig:proposed_system}の(3)）．
    \begin{itemize}
        \item \textbf{Hypo A, B:} ユーザーのフィードバックに忠実な改善案．
        \item \textbf{Hypo C:} LLMの音楽知識に基づき，今までと異なる表現解釈や意外性を含んだ提案（Deviated hypothesis）．
    \end{itemize}
    \item 3つの仮説をRenderBoxに入力し，3種類の候補演奏を生成する（図\ref{fig:proposed_system}の(4)(5)）．
    \item ユーザーは3つの候補演奏および基準演奏を聴き比べ，フィードバックを行うとともに，最も好ましい演奏を選択する．選択された演奏は次のループの「基準演奏」となる（図\ref{fig:proposed_system}の(6)(1)）．
    \item ユーザーが満足するか，ループ上限回数に達するまで繰り返す．
\end{enumerate}

\subsection{評価指標}
\label{subsec:評価指標}

実験結果の分析には，以下の観点および指標を用いる．

\subsubsection*{相互創発性 (Co-creativity)}
エージェント（AI）からの提案が，ユーザーの意図を超えて新しい価値を生み出したかを評価する．
\begin{itemize}
    \item \textbf{客観的評価:} 実験2において，LLMが独自の解釈を加えた「Music C（Hypo Cに基づく演奏）」が，次のループの基準演奏としてユーザーに選択された回数や割合．
    \item \textbf{主観的評価:} 「Music C」の演奏やその解釈テキストに対し，ユーザーが「意外だが良い」「新しい発見があった」と感じたかをアンケート（Q7, Q8）により評価する．
\end{itemize}

\subsubsection*{ユーザビリティ (Usability)}
システムが効率的かつ意図通りに機能したかを評価する．
\begin{itemize}
    \item \textbf{客観的評価:} ユーザーが最終的な満足に至るまでに要したループ回数を比較する．実験1より実験2の方が少ない回数で満足した場合，効率が良いと判断する．
    \item \textbf{主観的評価:} 「意図通りの演奏が生成されたか」「プロンプト生成の手間が軽減されたか」といった点をアンケート（Q4, Q6）により評価する．
\end{itemize}

実験に使用したアンケート項目の一覧を表\ref{tab:questionnaire_items}に示す．

\begin{table}[htbp]
\centering
\caption{被験者実験で使用したアンケート項目一覧}
\label{tab:questionnaire_items}
\resizebox{\textwidth}{!}{%
\begin{tabular}{llp{7cm}l}
\toprule
\textbf{Category} & \textbf{ID} & \textbf{Question Item} & \textbf{Scale} \\ \hline
\midrule
% --- 属性 ---
\multirow{4}{*}{Demographics} 
 & D1 & 年齢・性別 & Open \\
 & D2 & 普段音楽を聴く頻度・ジャンル & Open \\
 & D3 & 楽器演奏経験の有無とその年数 & Open \\
 & D4 & （事務用）被験者ID・名前 & - \\ \hline
\midrule
% --- 全体評価 ---
\multirow{4}{*}{\shortstack{Overall\\Satisfaction}} 
 & Q1 & セクション1（手動）で生成された演奏に満足しましたか？ & 5-point Likert \\
 & Q2 & セクション2（対話）で生成された演奏に満足しましたか？ & 5-point Likert \\
 & Q3 & どちらのセクションの演奏により満足しましたか？ & 2-AFC \\
 & Q4 & どちらのシステムが使いやすかったですか？ & 2-AFC \\ \hline
\midrule
% --- システム評価 ---
\multirow{2}{*}{\shortstack{System\\Capability}} 
 & Q5 & エージェントはあなたの理想の演奏を正しく推測していましたか？ & 5-point Likert \\
 & Q6 & あなたの指示通りに演奏が変化していましたか？ & 5-point Likert \\ 
\midrule \hline
% --- 相互創発性 (ここが重要) ---
\multirow{4}{*}{\shortstack{Co-creativity\\\& Interaction}} 
 & Q7 & エージェントの提案に，意図と異なりつつも「良い」と思えるものはありましたか？ & 5-point Likert \\
 & Q8 & 対話を通じて，当初イメージしていた曲の目標や雰囲気が変化しましたか？ & 5-point Likert \\
 & Q9 & エージェントとの対話に関する感想・気づき & Free Text \\
 & Q10 & 実験全体の感想 & Free Text \\ \hline
\bottomrule
\end{tabular}%
}
\vspace{1ex}
\small \\
Note: 5-point Likert (1: Strongly Disagree -- 5: Strongly Agree), 2-AFC (2-Alternative Forced Choice)
\end{table}